\documentclass[12pt]{article}
\usepackage{ceai}

\begin{document}

\title{CE 488/5588: Applied AI for Engineering \\ Final Assignment - Set B}
\author{}
\date{Spring 2026}

\maketitle

\begin{center}
\pmb{Name: \rule{3in}{0.15mm}}
\end{center}

\begin{ch-box}
This final assignment evaluates your understanding of AI foundations and machine learning methods in engineering. The problems are designed to progress from conceptual understanding to basic analytical formulation. Answer all parts of each problem.
\end{ch-box}

\section{Problem 1: Engineering Logic \& Data Representation}
\textbf{Context:} Homework 01-02 discussed the shift in engineering skills and data types.

\begin{enumerate}[label=(\alph*)]
    \item \textbf{Reskilling:} The WEF Future of Jobs report emphasizes the need for engineers to pivot from pure calculation to "AI-augmented design." Describe one engineering workflow that can be \textbf{augmented} by AI versus one that can be \textbf{automated}. Justify your classification for each. \\
    \fullunderline \\
    \fullunderline \\
    \fullunderline

    \item \textbf{Unsupervised vs Supervised:} You have a dataset of 5,000 photos of pavement. 
    \begin{itemize}
        \item If you use a model to group similar-looking cracks \textit{without} prior labels, is this Supervised or Unsupervised? Explain why. \rule{1.5in}{0.15mm} \\
        \fullunderline
        \item If you train a model to predict the "Pavement Condition Index (PCI)" from these photos using historical ratings, is this Supervised or Unsupervised? Explain why. \rule{1.5in}{0.15mm} \\
        \fullunderline
    \end{itemize}

    \item \textbf{PCA Logic:} When applying PCA to a dataset with 20 features, you find that the first 2 principal components explain 95\% of the variance. What does this tell you about the dimensionality of the underlying engineering problem? \\
    \fullunderline \\
    \fullunderline
\end{enumerate}

\section{Problem 2: Linear Models \& Slope Stability}
\textbf{Context:} Consider a binary classifier for slope stability ($c=1$ for stable, $c=0$ for failure) using $x_1$ (slope angle) and $x_2$ (cohesion).

\begin{enumerate}[label=(\alph*)]
    \item \textbf{Formulation:} A perceptron has weights $\bm{w} = [-3.0, 1.0]^T$ and bias $b = 5.0$. 
    \begin{itemize}
        \item Write the linear decision function $u(\bm{x})$. \rule{2.5in}{0.15mm} \\
        \fullunderline
        \item For a slope with $x_1 = 30^\circ$ and $x_2 = 100$ kPa, calculate the score $u$ and explain what its numerical value implies about the distance to the boundary. \rule{1in}{0.15mm} \\
        \fullunderline
    \end{itemize}

    \item \textbf{Decision Boundary:} Find the equation of the line where $u(\bm{x}) = 0$. Express it in the form $x_2 = m x_1 + c$. Show your derivation steps. \\
    \fullunderline \\
    \fullunderline

    \item \textbf{Linear Separability:} If your data contains two "Failure" points sandwiched between two "Stable" points along the $x_1$ axis, can a single Perceptron correctly classify all points? Why or why not? \\
    \fullunderline \\
    \fullunderline
\end{enumerate}

\section{Problem 3: Probabilistic Classification (Logistic Regression)}
\textbf{Context:} Logistic regression is widely used for risk assessment where a probability is required.

\begin{enumerate}[label=(\alph*)]
    \item A structural monitoring system predicts the probability of "Critical Failure" ($c=1$). For a given sensor input, the logit is $u = -2.2$. 
    \begin{itemize}
        \item Calculate the probability $P(c=1|\bm{x})$. (Use $e^{2.2} \approx 9.02$). \rule{1.5in}{0.15mm} \\
        \fullunderline
        \item If the critical threshold is set to $\tau = 0.1$ (due to high safety risk), what is the predicted state? Justify this threshold choice from an engineering safety perspective. \rule{1in}{0.15mm} \\
        \fullunderline
    \end{itemize}

    \item \textbf{Sensitivity:} Why might a structural engineer choose a threshold $\tau = 0.1$ instead of the default $0.5$? How does this affect the model's \textbf{Recall} and \textbf{False Alarm Rate}? Explain the trade-off. \\
    \fullunderline \\
    \fullunderline \\
    \fullunderline

    \item \textbf{Sigmoid Function:} Sketch the shape of the sigmoid function $\sigma(u)$. Label the axes and the points where $u=0$ and $u \to \pm \infty$. 
    \begin{ex-box}
    \vspace{2.5in}
    \end{ex-box}
\end{enumerate}

\section{Problem 4: Advanced Classification (SVM)}
\textbf{Context:} SVMs are powerful for high-dimensional engineering data.

\begin{enumerate}[label=(\alph*)]
    \item \textbf{Support Vectors:} In an SVM model, which data points are called "Support Vectors," and why are they the only points that determine the decision boundary? \\
    \fullunderline \\
    \fullunderline

    \item \textbf{Kernels:} Explain the conceptual difference between a \textbf{Linear Kernel} and a \textbf{Radial Basis Function (RBF) Kernel}. When would you choose the RBF kernel and why? \\
    \fullunderline \\
    \fullunderline \\
    \fullunderline

    \item \textbf{Slack Variables:} Engineering data is often "noisy" and not perfectly separable. How does the introduction of \textbf{Slack Variables} ($\xi$) allow the SVM to handle misclassified points in the training set? Explain the role of the penalty. \\
    \fullunderline \\
    \fullunderline
\end{enumerate}

\section{Problem 5: Neural Networks (MLP)}
\textbf{Context:} Deep learning models like MLPs are the building blocks of modern AI.

\begin{enumerate}[label=(\alph*)]
    \item \textbf{Architecture:} An MLP has:
    \begin{itemize}
        \item Input: 3 features.
        \item Hidden Layer 1: 5 neurons.
        \item Hidden Layer 2: 5 neurons.
        \item Output: 1 neuron.
    \end{itemize}
    Calculate the total number of \textbf{weights} (excluding biases) in this network. Show the count for each layer transition and explain your work. \\
    \fullunderline \\
    \fullunderline \\
    \fullunderline

    \item \textbf{Backpropagation:} Briefly describe the goal of the \textbf{Backpropagation} algorithm. What is being calculated (mathematically), and why is it essential for the learning process? \\
    \fullunderline \\
    \fullunderline

    \item \textbf{Vanishing Gradient:} Why is the \textbf{ReLU} activation function ($\max(0, u)$) often preferred over the Sigmoid function in very deep networks? Explain how it helps during training. \\
    \fullunderline \\
    \fullunderline
\end{enumerate}

\section{Problem 6: Modern AI Workflows}
\textbf{Context:} Topics 14-16 discussed the practical deployment of LLMs and CNNs.

\begin{enumerate}[label=(\alph*)]
    \item \textbf{CNN Properties:} In a convolutional layer, what is the effect of increasing the \textbf{Stride} from 1 to 2 on the size of the output feature map? Does it increase or decrease the "receptive field" overlap? Explain why. \\
    \fullunderline \\
    \fullunderline

    \item \textbf{Transformers:} Identify which Transformer architecture (\textbf{Encoder-only}, \textbf{Decoder-only}, or \textbf{Encoder-Decoder}) is most appropriate for the following tasks. Briefly justify each choice.
    \begin{itemize}
        \item Summarizing a 50-page geotechnical report: \rule{1.5in}{0.15mm} \\
        \fullunderline
        \item A chatbot that assists with Python coding for data analysis: \rule{1.5in}{0.15mm} \\
        \fullunderline
    \end{itemize}

    \item \textbf{Prompt Engineering:} You are using an LLM to extract "Material Type" and "Estimated Cost" from unstructured emails. Design a "Few-Shot" prompt template (briefly describe the components) to ensure the LLM returns structured JSON output. Explain why few-shot prompting is effective here. \\
    \fullunderline \\
    \fullunderline \\
    \fullunderline
\end{enumerate}


\end{document}
